\chapter{El \emph{matcher}, tier por tier}
\label{cap:matcher}

\begin{intuicion}
Llamar a \texttt{search\_temporal\_pattern} no ejecuta un solo
algoritmo --- ejecuta el operador físico que el \emph{planner} eligió
para esa hipótesis. Hay cuatro hoy, ordenados por especificidad y
\emph{lift} esperado: \emph{LFTJ-K4} (sólo cuando la hipótesis es un
\(K_4\) homogéneo), \emph{Yannakakis chain}, \emph{Yannakakis tree-walk}
(\emph{shape} \texttt{Graph}), y el \emph{DFS legacy} como
\emph{fallback}. Este capítulo explica cuándo se elige cada uno y por
qué.
\end{intuicion}

\section{El \emph{planner}}

\begin{lstlisting}[language=Rust]
pub enum DispatchPlan {
    DfsMatch,    // default, generic
    K4Lftj,      // M1/S3, sólo K_4 homogéneo
    Yannakakis,  // L1, α-acíclica + flag GRAPHHUNTER_YANNAKAKIS
}

pub fn plan(hypothesis: &Hypothesis) -> DispatchPlan {
    if name_marker_dispatch(hypothesis)   { return K4Lftj; }
    if structural_dispatch(hypothesis)    { return K4Lftj; }
    if yannakakis_dispatch(hypothesis)    { return Yannakakis; }
    DfsMatch
}
\end{lstlisting}

\begin{itemize}
    \item \texttt{name\_marker\_dispatch}: la \emph{back-compat} S3.
    Si el nombre de la hipótesis contiene la cadena ``k4-clique'' y
    \texttt{GRAPHHUNTER\_K4\_LFTJ=1}, ruta al LFTJ.
    \item \texttt{structural\_dispatch} (M1): detecta una hipótesis
    de exactamente 6 pasos homogéneos (mismo \texttt{EntityType},
    misma \texttt{RelationType}, sin predicados) sin requerir la
    palabra mágica.
    \item \texttt{yannakakis\_dispatch} (L1): cualquier hipótesis
    \(\alpha\)-acíclica con \texttt{GRAPHHUNTER\_YANNAKAKIS=1}.
    Ambas \emph{shapes} \texttt{Chain} y \texttt{Graph} aplican.
\end{itemize}

\section{LFTJ-K4: el \emph{worst-case-optimal join} para cliques}

Para una hipótesis \(K_4\) (cuatro vértices, las seis aristas
dirigidas entre todos los pares), el DFS recursivo escala como
\(O(m \cdot \deg^3)\) en el peor caso. \emph{Leapfrog Triejoin}
(Veldhuizen ICDT 2014) se queda en el \emph{AGM bound} \(m^{\rho^*} =
m^2\) — es \emph{worst-case optimal}.

\begin{intuicion}
La idea es construir un \emph{trie} sobre las aristas indexado por
\texttt{(source, dest)} y usar \emph{multi-way intersection} en lugar
de \emph{nested loop}. Para \(K_4\), una vez fijado un par
\((a, b)\), el conjunto de \(c\)s que extiende a un triángulo es
\(N(a) \cap N(b)\) (\emph{intersect SIMD}); el conjunto de \(d\)s que
completa el \(K_4\) es \(\{c\} \cap N(a) \cap N(b) \cap N(c)\).
\end{intuicion}

\begin{enelcodigo}
\filepath{core/graph-engine/src/lftj.rs} y
\filepath{core/graph-engine/src/lftj\_query.rs}. Versiones canónicas
(con \emph{symmetry breaking} hardcoded) en M3++ entregaron el corte
adicional \(|\mathrm{Aut}|\)-veces.
\end{enelcodigo}

\section{DFS legacy con \emph{prunes}}

Cuando ningún operador especializado dispatcha, el motor cae al DFS
recursivo de toda la vida. Pero el DFS de hoy lleva tres \emph{prunes}
acumulados:

\begin{enumerate}
    \item \textbf{NLF}: descarta \emph{starts} cuyo histograma de
    vecinos no satisface el primer paso.
    \item \textbf{k-hop reachability (D5/M2)}: descarta \emph{starts}
    desde los que no puede haber un camino de longitud
    \(\leq \mathrm{MAX\_DEPTH}\) al \texttt{dest\_type} final, dentro
    de la ventana temporal.
    \item \textbf{B9 \emph{symmetry breaking}} para hipótesis
    palindrómicas.
\end{enumerate}

Cada \emph{prune} se aplica antes del DFS principal. Capítulo~\ref{cap:indexes}
detalla cómo se construyen.

\section{Yannakakis chain (L1 stage 2)}

Para cualquier hipótesis \(\alpha\)-acíclica con \emph{shape}
\texttt{Chain}, el motor materializa los \emph{bags} y corre la
\emph{semi-join reducer} de dos pasadas antes de enumerar:

\begin{algorithm}[H]
\caption{Yannakakis chain matcher}
\KwIn{\texttt{bags}: \([B_0, B_1, \ldots, B_{n-1}]\)}
\KwOut{\texttt{Vec<HuntResult>}}
\tcc{Pasada 1 (bottom-up).}
\For{$i$ from $n-2$ \textbf{down to} $0$}{
    \texttt{live\_origins} $\leftarrow \{r.\mathrm{origin} : r \in B_{i+1}\}$\;
    $B_i \leftarrow \{r \in B_i : r.\mathrm{dest} \in \texttt{live\_origins}\}$\;
}
\tcc{Pasada 2 (top-down).}
\For{$i$ from $1$ \textbf{to} $n-1$}{
    \texttt{live\_dests} $\leftarrow \{r.\mathrm{dest} : r \in B_{i-1}\}$\;
    $B_i \leftarrow \{r \in B_i : r.\mathrm{origin} \in \texttt{live\_dests}\}$\;
}
\tcc{Enumerate.}
\For{cada \texttt{row} en $B_0$}{
    \texttt{walk}(\texttt{row.origin}, \texttt{row.dest},
                 \texttt{row.timestamp}, $1$)\;
}
\end{algorithm}

\begin{figure}[H]
\centering
\begin{tikzpicture}[
    node distance=0.3cm,
    every node/.style={font=\scriptsize},
    bag/.style={draw=accent, thick, rounded corners=1mm,
                minimum width=2.4cm, minimum height=1cm,
                fill=accentlight, align=center}
]
\node[bag] (b0) {$B_0$\\\itshape User--Auth--Host};
\node[bag, right=of b0] (b1) {$B_1$\\\itshape Host--Connect--IP};
\node[bag, right=of b1] (b2) {$B_2$\\\itshape IP--Beacon--Domain};

\draw[->, >=Latex, accent, thick, dashed] (b1.south) to[out=-90,in=-90, looseness=1.2] node[below, yshift=-0.2cm]{semi-join} (b0.south);
\draw[->, >=Latex, accent, thick, dashed] (b2.south) to[out=-90,in=-90, looseness=1.2] (b1.south);
\draw[->, >=Latex, accent!60, thick] (b0.north) to[out=90,in=90, looseness=1.2] node[above, yshift=0.2cm]{semi-join} (b1.north);
\draw[->, >=Latex, accent!60, thick] (b1.north) to[out=90,in=90, looseness=1.2] (b2.north);
\end{tikzpicture}
\caption{Las dos pasadas del \emph{semi-join reducer}. Bottom-up
elimina \emph{rows} cuyo \emph{dest} no aparece como \emph{origin} en
el siguiente \emph{bag}; top-down hace lo simétrico. Después de las
dos pasadas, toda \emph{row} sobreviviente participa de al menos un
\emph{full join}.}
\label{fig:matcher-semijoin}
\end{figure}

\section{Yannakakis tree-walk (L1 stage 3b)}

Para \emph{shape} \texttt{Graph} (estrellas, árboles, cualquier
\(\alpha\)-acíclico no-cadena), el \emph{enumerator} camina el
\emph{join tree} construido por GYO. Mantiene un
\texttt{binding: Vec<Option<StrId>>} indexado por \texttt{VarId}; en
cada nodo del árbol consulta el \emph{bag} correspondiente a través
del \emph{shared variable} con el padre.

\begin{algorithm}[H]
\caption{Tree-walk enumerator}
\KwIn{\texttt{join\_tree}, \texttt{bags}, \texttt{binding} (todo \texttt{None})}
\KwOut{\texttt{Vec<HuntResult>}}
\texttt{recurse}(\texttt{node\_idx} = 0)\;
\BlankLine
\Comment{\texttt{recurse(node\_idx)}}
\If{\texttt{node\_idx} == \texttt{join\_tree.nodes.len()}}{
    \texttt{emit}(\texttt{binding} resuelto a \texttt{Vec<Arc<str>>})\;
    \Return
}
\texttt{step} $\leftarrow$ \texttt{join\_tree.nodes[node\_idx]}\;
\texttt{edge} $\leftarrow$ \texttt{hypergraph.edges[step]}\;
\Switch{\texttt{(binding[edge.origin], binding[edge.dest])}}{
    \Case{\texttt{(Some(o), \_)}}{candidates $\leftarrow$ \texttt{by\_origin[step][o]}}
    \Case{\texttt{(None, Some(d))}}{candidates $\leftarrow$ \texttt{by\_dest[step][d]}}
    \Case{\texttt{(None, None)}}{candidates $\leftarrow$ \texttt{bags[step].rows} (root)}
}
\For{\texttt{row} \textbf{in} candidates}{
    intentar \texttt{bind(edge.origin, row.origin)} y \texttt{bind(edge.dest, row.dest)}\;
    \texttt{recurse(node\_idx + 1)}\;
    \texttt{unbind} si correspondía\;
}
\end{algorithm}

\begin{intuicion}
La diferencia clave con el \emph{chain enumerator}: en una cadena
sabemos que \texttt{bag[i+1].origin == bag[i].dest} por la convención
posicional, así que el \emph{lookup} es siempre \emph{by\_origin}. En
una estrella, el \emph{shared variable} con el padre puede aparecer
en cualquiera de los dos lados; por eso el \emph{tree-walk} mantiene
ambos índices y dispatcha por \texttt{(origin\_bound, dest\_bound)}.
\end{intuicion}
